\documentclass[10pt]{beamer}

% Tema Limpo e Institucional IFCE Taua
\usetheme{Madrid}
\usecolortheme{whale}

% Pacotes Essenciais
\usepackage[utf8]{inputenc}
\usepackage[T1]{fontenc}
\usepackage[brazilian]{babel}
\usepackage{graphicx}
\usepackage{booktabs}
\usepackage{amsmath}
\usepackage{amssymb}
\usepackage{listings}
\usepackage{xcolor}
\usepackage{tikz}
\usetikzlibrary{shapes.geometric,arrows.meta,positioning,calc}

% Configuracao de Cores e Codigo (Python)
\definecolor{codegreen}{rgb}{0,0.6,0}
\definecolor{codegray}{rgb}{0.5,0.5,0.5}
\definecolor{codepurple}{rgb}{0.58,0,0.82}
\definecolor{backcolour}{rgb}{0.96,0.97,0.98}

\lstset{
    language=Python,
    backgroundcolor=\color{backcolour},
    basicstyle=\ttfamily\scriptsize,
    keywordstyle=\color{blue}\bfseries,
    stringstyle=\color{red!80!black},
    commentstyle=\color{green!60!black},
    numberstyle=\tiny\color{codegray},
    breaklines=true,
    showstringspaces=false,
    tabsize=4,
    frame=single,
    rulecolor=\color{blue!20}
}

% Metadados da Apresentacao
\title[Engenharia de Software -- Aula 26]{Engenharia de Software -- Aula 26}
\subtitle{Sprint 04: Diagrama de Sequência e Orquestração do Negócio}
\author[Prof. Me. Reginaldo Fernandes]{Prof. Me. Reginaldo Pereira Fernandes}
\institute[IFCE Campus Tauá]{
    Instituto Federal de Educação, Ciência e Tecnologia do Ceará -- Campus Tauá\\
    Técnico em Informática para Internet (TSII.210)
}
\date{Semana 2 -- Encontro 3 (100 min)}

\begin{document}

% Slide 1: Titulo
\begin{frame}
    \titlepage
\end{frame}

% Slide 2: Agenda do Encontro
\begin{frame}{Agenda do Encontro 3 (100 Minutos)}
    \tableofcontents
\end{frame}

\section{Modelagem Dinâmica: Diagrama de Sequência}

\begin{frame}{Por que Precisamos da Visão Dinâmica?}
    \begin{itemize}
        \item \textbf{O que já temos até a Aula 25:}
        \begin{itemize}
            \item Entidade \texttt{Cupom} com regras de negócio em Python.
            \item Tabela \texttt{cupons} e \texttt{CupomRepository} persistindo no SQLite.
        \end{itemize}
        \item \textbf{A Pergunta que o Diagrama de Classes Não Responde:}
        \begin{itemize}
            \item \textit{"Em que ordem temporal os objetos conversam durante a compra de um livro com cupom?"}
            \item \textit{"Quem deve consultar o repositório? Quem deve debitar o estoque? Quem deve atualizar o total?"}
        \end{itemize}
        \item \textbf{Diagrama de Sequência (UML 2 -- Cap. 5):}
        \begin{itemize}
            \item O padrão industrial para modelar o fluxo de execução síncrono e transacional entre componentes.
        \end{itemize}
    \end{itemize}
\end{frame}

\begin{frame}{Diagrama de Sequência: Fluxo do Checkout com Cupom}
\centering
\resizebox{0.95\textwidth}{!}{%
\begin{tikzpicture}[>=Latex, font=\sffamily]
    % Participantes
    \node[draw=blue!80!black, fill=blue!10, rounded corners, font=\scriptsize\bfseries] (cli) at (0, 0) {Cliente (UI)};
    \node[draw=green!60!black, fill=green!10, rounded corners, font=\scriptsize\bfseries] (srv) at (3.2, 0) {CheckoutService};
    \node[draw=yellow!80!black, fill=yellow!15, rounded corners, font=\scriptsize\bfseries] (rep) at (6.4, 0) {CupomRepo};
    \node[draw=purple!80!black, fill=purple!10, rounded corners, font=\scriptsize\bfseries] (ped) at (9.6, 0) {Pedido};

    % Linhas de vida
    \draw[dashed, gray] (cli) -- ++(0, -4.6) coordinate (cli_end);
    \draw[dashed, gray] (srv) -- ++(0, -4.6) coordinate (srv_end);
    \draw[dashed, gray] (rep) -- ++(0, -4.6) coordinate (rep_end);
    \draw[dashed, gray] (ped) -- ++(0, -4.6) coordinate (ped_end);

    % Mensagens com timestamps verticais
    \draw[->, thick, draw=blue!80!black] ($(cli)!0.12!(cli_end)$) -- node[above, font=\tiny\bfseries] {1: checkout(pedido, COD)} ($(srv)!0.12!(srv_end)$);
    \draw[->, thick, draw=black!80] ($(srv)!0.26!(srv_end)$) -- node[above, font=\tiny\bfseries] {2: buscar\_por\_codigo(COD)} ($(rep)!0.26!(rep_end)$);
    \draw[<-, dashed, thick, draw=black!60] ($(srv)!0.38!(srv_end)$) -- node[above, font=\tiny] {3: cupom} ($(rep)!0.38!(rep_end)$);
    \draw[->, thick, draw=black!80] ($(srv)!0.52!(srv_end)$) -- node[above, font=\tiny\bfseries] {4: aplicar\_cupom(cupom)} ($(ped)!0.52!(ped_end)$);
    \draw[<-, dashed, thick, draw=black!60] ($(srv)!0.64!(srv_end)$) -- node[above, font=\tiny] {5: total líquido} ($(ped)!0.64!(ped_end)$);
    \draw[->, thick, draw=orange!90!black] ($(srv)!0.76!(srv_end)$) -- node[above, font=\tiny\bfseries] {6: registrar\_uso(id)} ($(rep)!0.76!(rep_end)$);
    \draw[<-, dashed, thick, draw=blue!80!black] ($(cli)!0.90!(cli_end)$) -- node[above, font=\tiny\bfseries] {7: Pedido Confirmado} ($(srv)!0.90!(srv_end)$);
\end{tikzpicture}
}
\end{frame}

\section{Diagrama de Atividades e Regras de Negócio}

\begin{frame}{Diagrama de Atividades: Ciclo de Decisão no Checkout}
\centering
\resizebox{!}{0.72\textheight}{%
\begin{tikzpicture}[>=Latex, font=\sffamily, node distance=0.8cm]
    \node[draw, circle, fill=black, minimum size=0.35cm] (start) {};
    \node[draw=blue!80!black, fill=blue!5, rounded corners, below=0.4cm of start, font=\scriptsize] (inicia) {Recebe Pedido e Código de Cupom};
    \node[draw=purple!80!black, fill=purple!5, diamond, aspect=1.8, below=0.5cm of inicia, font=\tiny\bfseries, align=center] (dec_cupom) {Cupom\\informado?};
    \node[draw=green!60!black, fill=green!5, rounded corners, below=0.6cm of dec_cupom, font=\scriptsize] (valida) {Valida Regras de Elegibilidade (DUC)};
    \node[draw=green!60!black, fill=green!5, rounded corners, below=0.5cm of valida, font=\scriptsize] (aplica) {Abate Desconto e Calcula Total Líquido};
    \node[draw=orange!80!black, fill=orange!5, rounded corners, below=0.5cm of aplica, font=\scriptsize] (baixa) {Debita Estoque do Livro e Incrementa Uso};
    \node[draw, circle, double, fill=black, minimum size=0.4cm, below=0.4cm of baixa] (stop) {};

    \draw[->, thick] (start) -- (inicia);
    \draw[->, thick] (inicia) -- (dec_cupom);
    \draw[->, thick] (dec_cupom) -- node[midway, right, font=\tiny\bfseries] {Sim} (valida);
    \draw[->, thick] (dec_cupom.east) -- ++(1.5,0) |- node[near end, right, font=\tiny\bfseries] {Não (Sem Cupom)} (baixa.east);
    \draw[->, thick] (valida) -- (aplica);
    \draw[->, thick] (aplica) -- (baixa);
    \draw[->, thick] (baixa) -- (stop);
\end{tikzpicture}
}
\end{frame}

\section{Implementação: Calculadora de Desconto}

\begin{frame}[fragile]{Mão na Massa: \texttt{CalculadoraDescontoService}}
\begin{lstlisting}[language=Python]
class CalculadoraDescontoService:
    """Servico de dominio puro: calculo matematico de descontos."""

    @staticmethod
    def calcular(subtotal: float, tipo_desconto: TipoDesconto, valor_desconto: float) -> float:
        if subtotal <= 0:
            return 0.0

        if tipo_desconto == TipoDesconto.PERCENTUAL:
            desconto = subtotal * (valor_desconto / 100.0)
        elif tipo_desconto == TipoDesconto.VALOR_FIXO:
            desconto = min(valor_desconto, subtotal)
        else:
            raise ValueError(f"Tipo de desconto desconhecido: {tipo_desconto}")

        return round(desconto, 2)
\end{lstlisting}
    \begin{block}{Por que um Serviço Isolado?}
    Permite testar e reutilizar as fórmulas matemáticas de desconto sem depender de banco de dados ou da interface!
    \end{block}
\end{frame}

\section{Implementação: Integração com o Pedido}

\begin{frame}[fragile]{Mão na Massa: Integrando o Cupom ao \texttt{Pedido}}
\begin{lstlisting}[language=Python]
class Pedido:
    def __init__(self, cliente: Cliente, id: Optional[int] = None):
        self.id = id
        self.cliente = cliente
        self.itens: List[ItemPedido] = []
        self.cupom: Optional[Cupom] = None
        self.desconto_aplicado: float = 0.0
        self.status: str = "CRIADO"

    @property
    def subtotal(self) -> float:
        return round(sum(item.subtotal for item in self.itens), 2)

    def aplicar_cupom(self, cupom: Cupom) -> float:
        desconto = cupom.calcular_desconto(self.subtotal)
        self.cupom = cupom
        self.desconto_aplicado = desconto
        return self.desconto_aplicado

    @property
    def total(self) -> float:
        return max(round(self.subtotal - self.desconto_aplicado, 2), 0.0)
\end{lstlisting}
\end{frame}

\section{Implementação: CheckoutService Orquestrador}

\begin{frame}[fragile]{Mão na Massa: \texttt{CheckoutService} Transacional}
\begin{lstlisting}[language=Python]
class CheckoutService:
    def __init__(self, cupom_repo: CupomRepository, livro_repo: LivroRepository, pedido_repo: PedidoRepository):
        self.cupom_repo = cupom_repo
        self.livro_repo = livro_repo
        self.pedido_repo = pedido_repo

    def processar_checkout(self, pedido: Pedido, codigo_cupom: Optional[str] = None) -> Pedido:
        # 1. Se informou cupom, busca e aplica
        if codigo_cupom:
            cupom = self.cupom_repo.buscar_por_codigo(codigo_cupom)
            if not cupom:
                raise CupomInexistenteError(codigo_cupom)
            pedido.aplicar_cupom(cupom)

        # 2. Debita estoque dos livros adquiridos
        for item in pedido.itens:
            self.livro_repo.atualizar_estoque(item.livro.id, -item.quantidade)

        # 3. Registra uso do cupom se aplicado
        if pedido.cupom:
            self.cupom_repo.registrar_uso(pedido.cupom.id)

        # 4. Persiste o pedido e finaliza status
        pedido.status = "PAGO"
        return self.pedido_repo.salvar(pedido)
\end{lstlisting}
\end{frame}

\section{Rito Ágil e Próximos Passos}

\begin{frame}{Rito Ágil: Registro da Daily 3}
    \begin{block}{Anotação para o Relatório da Sprint no \texttt{guia\_desenvolvedor.md}}
    \begin{itemize}
        \item \textbf{O que realizei hoje?}
        \begin{itemize}
            \item Interpretei o Diagrama de Sequência UML do fluxo de checkout.
            \item Implementei a \texttt{CalculadoraDescontoService}.
            \item Integrei o agregado \texttt{Pedido} ao cupom com recálculo automático do total líquido.
            \item Implementei o orquestrador \texttt{CheckoutService} garantindo baixa de estoque e registro de uso.
        \end{itemize}
        \item \textbf{O que farei no Encontro 4 (Aula 27)?}
        \begin{itemize}
            \item Construir a suíte de testes unitários e de integração com \textbf{Pytest}.
            \item Demonstrar o fluxo interativo no CLI e submeter o Pull Request final.
        \end{itemize}
        \item \textbf{Impedimentos?} Nenhum impedimento.
    \end{itemize}
    \end{block}
\end{frame}

\begin{frame}{Síntese do Encontro 3}
    \begin{itemize}
        \item \textbf{Sistema Totalmente Integrado:}
        \begin{itemize}
            \item A modelagem dinâmica se materializou em serviços coordenados.
            \item O cálculo de desconto e a baixa de estoque funcionam de ponta a ponta.
        \end{itemize}
        \item \textbf{Para o Encontro 4 (Aula 27 - Fechamento da Sprint):}
        \begin{itemize}
            \item Validação de Qualidade: \textbf{Pytest} 100\% verde.
            \item Demonstração CLI para a liderança.
            \item Submissão do Pull Request e Retrospectiva da Sprint.
        \end{itemize}
    \end{itemize}
    \vspace{0.4cm}
    \centering
    \textbf{Excelente trabalho! O Módulo de Cupons está quase pronto para entrega!}
\end{frame}

\end{document}
